\documentclass[12pt]{article}
\usepackage[utf8]{inputenc}
\usepackage[spanish]{babel}
\usepackage{amsmath}
\usepackage{amsfonts}
\usepackage{amssymb}
\usepackage{geometry}
\geometry{letterpaper, margin=2.5cm}

\begin{document}

\begin{center}
    \Large \textbf{Evaluación de Examen 1: Unidad I} \\
    \large Física para Ciencias de la Salud (005-1713) \\
    \vspace{0.5cm}
    \textbf{Estudiante:} Daniela Romero \quad \textbf{C.I.:} 31.435.612
    \vspace{0.3cm} \\
    \large \textbf{Calificación Final: 8/10 puntos}
\end{center}

\vspace{0.5cm}

\section*{Análisis de Resultados}

\subsection*{1. Ángulo plano (Conversión a radianes)}
\textbf{Procedimiento y resultado:} Plantea de forma correcta la conversión multiplicando por el factor $\frac{\pi}{180^\circ}$. Simplifica la fracción con agilidad dividiendo entre 15 tanto el numerador como el denominador, obteniendo directamente el valor exacto de $\frac{5\pi}{12}$ rad. \\
\textbf{Veredicto:} \textbf{Correcto} (2/2 pts).

\subsection*{2. Sistemas de coordenadas (Longitud del hueso)}
\textbf{Procedimiento y resultado:} Calcula correctamente las distancias en cada eje restando las coordenadas de los extremos ($a = 7-1=6$ y $b = 10-2=8$). Luego aplica el Teorema de Pitágoras ($\sqrt{6^2 + 8^2} = \sqrt{100}$) llegando sin errores al resultado correcto de $10$ cm. \\
\textbf{Veredicto:} \textbf{Correcto} (2/2 pts).

\subsection*{3. Vectores (Fuerza resultante)}
\textbf{Procedimiento y resultado:} Plantea la suma de las componentes $\hat{i}$ y $\hat{j}$ de manera directa, respetando las leyes de los signos y obteniendo el vector exacto $30\hat{i} + 60\hat{j}$ N. Cabe mencionar como único detalle formal que denota el vector como $\Delta\vec{F}$ en lugar de $\vec{F}_R$, pero el cálculo aritmético de la sumatoria es impecable. \\
\textbf{Veredicto:} \textbf{Correcto} (2/2 pts).

\subsection*{4. Potencias de diez (Notación científica y prefijo)}
\textbf{Procedimiento y resultado:} La estudiante comete un error matemático al desplazar la coma decimal: expresa la cifra como $1,25 \times 10^{-6}$ m, lo cual equivaldría a $0,00000125$ m. La notación matemáticamente correcta debió ser $1,25 \times 10^{-5}$ m o en su defecto $12,5 \times 10^{-6}$ m. Adicionalmente, se omitió por completo la segunda parte de la pregunta, al no identificar ni nombrar el prefijo del S.I. \\
\textbf{Veredicto:} \textbf{Incorrecto} (0/2 pts).

\subsection*{5. Sistemas de unidades (Conversión de densidad)}
\textbf{Procedimiento y resultado:} Excelente dominio de los factores de conversión en cadena. Expresa y aplica correctamente las relaciones $1 \text{ kg} = 1000 \text{ g}$ y $1 \text{ m}^3 = 1.000.000 \text{ cm}^3$. Desarrolla las multiplicaciones y divisiones paso a paso de forma muy clara, llegando al resultado exacto de $1030 \text{ kg/m}^3$. \\
\textbf{Veredicto:} \textbf{Correcto} (2/2 pts).

\vspace{0.5cm}
\section*{Comentarios Generales y Sugerencias}
La estudiante demuestra un buen nivel general en su razonamiento analítico, con excelentes procedimientos algebraicos y un dominio muy destacable en la aplicación de los factores de conversión en cadena.
\begin{itemize}
    \item Se sugiere repasar cuidadosamente las reglas para expresar números en notación científica, asegurándose de que el exponente de la base 10 corresponda exactamente a la cantidad de lugares que se desplaza la coma decimal.
    \item Se recomienda leer los enunciados en su totalidad para no omitir partes esenciales de la respuesta, como sucedió en la Pregunta 4.
\end{itemize}

\end{document}